\documentclass[11pt]{article}
\usepackage{amsmath,amssymb,amsfonts}
\usepackage{hyperref}
\usepackage{graphicx}
\usepackage{geometry}
\geometry{margin=1in}

\title{A Dirac-Type Field Model for Primes}
\author{Thomas Hoffbauer}
\date{\today}

\begin{document}
\maketitle

\begin{abstract}
We propose a Dirac-type field model on the number line in which the
non-trivial zeros of the Riemann zeta function act as spectral modes,
prime numbers as negative charges, and prime gaps as positive artifacts.
A smooth-cut construction $W(p)$ with a mirror map defines local
particle--antiparticle pairs. Numerical experiments based on the first
$1000$ Riemann zeros show (i) GUE statistics of the spectrum, (ii) phase
antisymmetry of mirror pairs, and (iii) resonance of large prime gaps
with spectral modes, suggesting a spectral--geometric duality between
zeros and primes.
\end{abstract}

\section{Introduction}
The connection between prime numbers and the zeros of the Riemann zeta
function has motivated operator-theoretic interpretations since
Hilbert--P\'olya. In this work we develop a Dirac-type field model in
which primes appear as discrete charges interacting with a spectral
field defined by the zeta zeros.

\section{Smooth Cuts and Mirror Primes}
Let $H_3=\{2^a3^b\}$ and $H_5=\{2^a3^b5^c\}$. For each prime $p$ define
\begin{equation}
W(p)=\big(\max(\ell_3(p),\ell_5(p)),\min(u_3(p),u_5(p))\big),
\end{equation}
where $\ell_s(p)$ and $u_s(p)$ denote the neighboring $s$-smooth numbers.
Normalized coordinates $\theta_s$ yield a mirror map
$\theta_s^*(p)=1-\theta_s(p)$ and a mirror prime $p^\complement$ chosen by
least-square matching in the $(\theta_3,\theta_5)$ plane.

\section{Dirac-Type Field}
Define the field
\begin{equation}
\Phi(x)=\sum_{n\le N}\Re\!\left(\frac{x^{1/2+i\gamma_n}}{1/2+i\gamma_n}\right),
\end{equation}
where $\rho_n=1/2+i\gamma_n$ are the non-trivial zeros of $\zeta(s)$. The
associated phase field is
\begin{equation}
S(x)=\sum_{n\le N} e^{i\gamma_n\log x},\qquad \varphi(x)=\arg S(x).
\end{equation}

\section{Numerical Experiments}
Using the first $1000$ Riemann zeros we perform three tests:
\begin{enumerate}
\item \textbf{Spectral statistics:} Nearest-neighbor spacings of the
zeros follow the GUE distribution (Wigner surmise).
\item \textbf{Mirror symmetry:} For mirror pairs $(p,p^\complement)$ we
observe approximate phase antisymmetry
$\varphi(p)\approx-\varphi(p^\complement)$ modulo $2\pi$.
\item \textbf{Resonance of gaps:} Large prime gaps tend to occur when
$\log m_k$ (midpoints of gaps) is close to a spectral period
$2\pi/\gamma_n$.
\end{enumerate}

\section{Discussion}
These findings support the interpretation of primes and gaps as charges
and artifacts in a quantized number-theoretic field whose stability axis
is the critical line $\Re(s)=1/2$.

\section{Outlook}
Future work includes extension to larger datasets and reconstruction of
the underlying operator from its spectral data.

\bibliographystyle{plain}
\begin{thebibliography}{9}
\bibitem{Odlyzko}
A.~M. Odlyzko, \emph{The $10^{20}$th zero of the Riemann zeta function and
40 million of its neighbors}, preprint.
\bibitem{Montgomery}
H.~L. Montgomery, \emph{The pair correlation of zeros of the zeta
function}, Proc. Symp. Pure Math. \textbf{24} (1973).
\end{thebibliography}

\end{document}



\section{Symmetry Groups, Cayley Tables, and Operator Algebra}

\subsection{From Classes to Permutations}
We label the four residue families of the Bamberg model by
$\mathcal{C}=\{e,a,b,c\}$ (corresponding to the modulo--12 families
$1,5,7,11$). Any transformation acting on these families induces a
permutation of $\mathcal{C}$. Hence the symmetry algebra of class
transformations embeds naturally into the symmetric group $S_4$.

\subsection{Group Elements and Cayley Table}
Let $\mathcal{G}$ denote the group generated by the following elementary
operations:
\begin{itemize}
\item $\sigma_1$: swap $e\leftrightarrow a$,
\item $\sigma_2$: swap $b\leftrightarrow c$,
\item $\kappa$: class complement (duality map),
\item $r$: cyclic shift $(e,a,b,c)\mapsto(a,b,c,e)$.
\end{itemize}
These generators define a subgroup $\mathcal{G}\subset S_4$.
For $g_i,g_j\in\mathcal{G}$ the Cayley law reads
\begin{equation}
g_i\circ g_j = g_k
\qquad\Longleftrightarrow\qquad
U(g_i)U(g_j)=U(g_k),
\end{equation}
where $U(g)$ is the permutation matrix associated with $g$.

\subsection{Matrix Representation}
Each $g\in\mathcal{G}$ is represented by a $4\times4$ permutation matrix
$U(g)$ acting on the class state space $\mathbb{C}^4$:
\begin{equation}
U(g)_{ij}=\delta_{i,g(j)}.
\end{equation}
The Cayley table of $\mathcal{G}$ therefore induces directly the operator
multiplication table of the matrices $U(g)$.

\section{Dirac Operator on Primes and Dual Gap Operator}

\subsection{Prime States and Gap States}
We define a Hilbert space $\mathcal{H}_P=\ell^2(\mathbb{N})$ with orthonormal
basis $\{|p_n\rangle\}$ labeled by primes $p_n$, and a dual space
$\mathcal{H}_G$ with basis $\{|g_n\rangle\}$ labeled by gaps
$g_n=p_{n+1}-p_n$.

\subsection{Dirac-Type Operator}
The Dirac operator on primes is defined spectrally by
\begin{equation}
D_P |p_n\rangle = p_n |p_n\rangle.
\end{equation}
Its dual on gaps reads
\begin{equation}
D_G |g_n\rangle = g_n |g_n\rangle.
\end{equation}
This establishes the fundamental duality:
\begin{center}
\emph{Dirac operator $=$ primes,\quad dual operator $=$ gaps.}
\end{center}

\subsection{Symmetry Action on Operators}
For every $g\in\mathcal{G}$ we define a unitary action
\begin{equation}
U(g)|p_n\rangle = |p_{g\cdot n}\rangle,
\end{equation}
inducing the conjugation rule
\begin{equation}
D_P \mapsto U(g)D_PU(g)^{-1}.
\end{equation}
If $[D_P,U(g)]=0$, the corresponding transformation is a symmetry of the
prime Dirac spectrum.

\section{Group-Theoretic Decomposition and Invariants}

\subsection{Subgroups and Sectors}
Following the paradigm of the tetrahedral group $T_d\simeq S_4$, the
symmetry group $\mathcal{G}$ decomposes into subgroups corresponding to
distinct structural scales:
\begin{itemize}
\item cyclic sectors (shift symmetries),
\item dihedral sectors (class swaps and reflections),
\item alternating sector $A_4$ (orientation-preserving transformations).
\end{itemize}
Each sector defines an invariant subspace of $\mathcal{H}_P$.

\subsection{Projectors}
For a subgroup $H\subset\mathcal{G}$ we define the projector
\begin{equation}
\Pi_H = \frac{1}{|H|}\sum_{h\in H} U(h).
\end{equation}
These projectors isolate symmetry-adapted components of the prime field.

\section{Interpretation in the Bamberg Model}

The Bamberg model gains a natural algebraic formulation:
\begin{itemize}
\item residue families $\leftrightarrow$ basis states,
\item Hoffbauer cuts $\leftrightarrow$ operators,
\item class complements $\leftrightarrow$ reflections,
\item smoothness scales $\leftrightarrow$ subgroup strata.
\end{itemize}
The tetrahedral paradigm shows that arithmetic structure can be organized
exactly as geometric symmetry: via groups, Cayley tables, and operator
representations.

\section{Conclusion}
By identifying the Dirac operator with the prime spectrum and its dual with
the gap spectrum, and by encoding arithmetic transformations through
group-theoretic and matrix-based methods inspired by the tetrahedral
group $T_d$ and $S_4$, we obtain a unified algebraic framework for the
Bamberg model and a Dirac-type theory of primes.



% =======================
% ARXIV-MATHEMATICAL EXTENSIONS
% =======================

\section{Explicit Group Structure and Matrix Realization}

\subsection{Generators and Relations}
We fix the four Bamberg classes $\mathcal C=\{e,a,b,c\}$ and identify them
with $\{1,2,3,4\}$. Define generators in $S_4$ by
\begin{align}
r &= (1234),\\
s &= (12),\\
t &= (34),\\
\kappa &= (13)(24).
\end{align}
They generate a subgroup $\mathcal G\subset S_4$ containing cyclic,
dihedral and alternating components. The defining relations are
\begin{equation}
r^4=e,\qquad s^2=t^2=\kappa^2=e,\qquad srs=r^{-1},\qquad [s,t]=e.
\end{equation}

\subsection{Cayley Table (Abstract Law)}
For generators $r,s$ of the dihedral core $D_4\subset\mathcal G$ the Cayley
table reads
\begin{center}
\begin{tabular}{c|cccccccc}
$\circ$ & $e$ & $r$ & $r^2$ & $r^3$ & $s$ & $sr$ & $sr^2$ & $sr^3$\\ \hline
$e$     & $e$ & $r$ & $r^2$ & $r^3$ & $s$ & $sr$ & $sr^2$ & $sr^3$\\
$r$     & $r$ & $r^2$ & $r^3$ & $e$ & $sr^3$ & $s$ & $sr$ & $sr^2$\\
$r^2$   & $r^2$ & $r^3$ & $e$ & $r$ & $sr^2$ & $sr^3$ & $s$ & $sr$\\
$r^3$   & $r^3$ & $e$ & $r$ & $r^2$ & $sr$ & $sr^2$ & $sr^3$ & $s$\\
$s$     & $s$ & $sr$ & $sr^2$ & $sr^3$ & $e$ & $r$ & $r^2$ & $r^3$\\
$sr$    & $sr$ & $sr^2$ & $sr^3$ & $s$ & $r^3$ & $e$ & $r$ & $r^2$\\
$sr^2$  & $sr^2$ & $sr^3$ & $s$ & $sr$ & $r^2$ & $r^3$ & $e$ & $r$\\
$sr^3$  & $sr^3$ & $s$ & $sr$ & $sr^2$ & $r$ & $r^2$ & $r^3$ & $e$
\end{tabular}
\end{center}

\subsection{Permutation Matrices}
Let $\{e_1,e_2,e_3,e_4\}$ denote the canonical basis of $\mathbb C^4$.
For $\pi\in S_4$ define
\begin{equation}
U(\pi)e_j=e_{\pi(j)},\qquad
U(\pi)_{ij}=\delta_{i,\pi(j)}.
\end{equation}
Explicitly, the generators are represented by
\begin{align}
U(r)&=
\begin{pmatrix}
0&0&0&1\\
1&0&0&0\\
0&1&0&0\\
0&0&1&0
\end{pmatrix},
&
U(s)&=
\begin{pmatrix}
0&1&0&0\\
1&0&0&0\\
0&0&1&0\\
0&0&0&1
\end{pmatrix},\\[1ex]
U(t)&=
\begin{pmatrix}
1&0&0&0\\
0&1&0&0\\
0&0&0&1\\
0&0&1&0
\end{pmatrix},
&
U(\kappa)&=
\begin{pmatrix}
0&0&1&0\\
0&0&0&1\\
1&0&0&0\\
0&1&0&0
\end{pmatrix}.
\end{align}
These matrices are orthogonal and unitary, hence define a finite group
representation $U:\mathcal G\to U(4)$.

\section{Spectral Symmetry of the Prime Dirac Operator}

\subsection{Commutation Relations}
For $g\in\mathcal G$ the symmetry condition reads
\begin{equation}
[D_P,U(g)]=0.
\end{equation}
In this case the eigenspaces of $D_P$ decompose into irreducible
$\mathcal G$-modules. If the commutator is non-zero, the corresponding
transformation represents a controlled symmetry breaking between prime
classes.

\subsection{Gap Duality}
The dual operator $D_G$ satisfies
\begin{equation}
D_G = \Delta D_P \Delta^{-1},
\end{equation}
where $\Delta|p_n\rangle=|g_n\rangle$ is the prime--gap correspondence map.
Hence prime and gap spectra are unitarily equivalent up to this duality
transform.

\section{Representation-Theoretic Decomposition}

\subsection{Isotypic Components}
Let $\widehat{\mathcal G}$ denote the set of irreducible representations of
$\mathcal G$. Then
\begin{equation}
\mathcal H_P = \bigoplus_{\rho\in\widehat{\mathcal G}} 
\mathcal H_P^{(\rho)},
\end{equation}
where $\mathcal H_P^{(\rho)}$ is the isotypic component corresponding to
$\rho$.

\subsection{Character Projectors}
For each irreducible character $\chi_\rho$ we define the projector
\begin{equation}
\Pi_\rho=\frac{d_\rho}{|\mathcal G|}\sum_{g\in\mathcal G}
\overline{\chi_\rho(g)}\,U(g).
\end{equation}
These projectors isolate the symmetry-adapted parts of the prime field and
yield a canonical decomposition of the Dirac spectrum.

\section{Mathematical Positioning}

The present framework places the Bamberg model and the Dirac-type theory
of primes in the context of
\begin{itemize}
\item finite group representations,
\item operator algebras generated by permutation symmetries,
\item spectral theory on discrete arithmetic spaces.
\end{itemize}
This positioning aligns the model with modern mathematical physics and
noncommutative geometry approaches to number theory.



% =======================
% ARXIV ABSTRACT & INTRO
% =======================

\begin{abstract}
We develop a Dirac-type operator framework for prime numbers in which the
prime sequence appears as the spectrum of a self-adjoint operator and the
prime gaps arise as the spectrum of a canonically dual operator. Building
on the Bamberg model of residue-class families, we introduce a finite
symmetry group acting on prime states, identified with a subgroup of
$S_4$ inspired by the tetrahedral symmetry paradigm. This leads to an
explicit operator algebra generated by permutation matrices, equipped
with Cayley tables and representation-theoretic projectors. The approach
places prime dynamics in the setting of finite group representations,
spectral theory, and noncommutative geometry, providing a unified
algebraic framework for primes and gaps.
\end{abstract}

\section{Introduction}

The idea that prime numbers should be understood spectrally has appeared
repeatedly in modern number theory and mathematical physics, most notably
in the Hilbert--P\'olya conjecture and in the Berry--Keating program. In
these approaches one seeks a self-adjoint operator whose eigenvalues
encode the non-trivial zeros of the Riemann zeta function or, more
generally, arithmetic spectra.

In this work we pursue a complementary perspective: instead of encoding
the zeros, we encode the primes themselves as the spectrum of a Dirac-type
operator. The guiding principle is the following duality:
\begin{center}
\emph{Primes form the spectrum of a Dirac operator, while gaps form the
spectrum of its canonical dual.}
\end{center}

The mathematical structure underlying this duality is provided by the
Bamberg model, which organizes primes into four residue-class families
modulo $12$. These families naturally support a finite symmetry group of
transformations. Inspired by the classical identification of the full
tetrahedral group $T_d$ with the symmetric group $S_4$, we model the
symmetry of prime families by a subgroup of $S_4$ acting on the class
labels. This step allows us to replace heuristic transformations by a
precise operator algebra generated by permutation matrices.

The main contributions of this paper are:
\begin{itemize}
\item the definition of a Dirac-type operator whose spectrum consists of
prime numbers;
\item the construction of a dual operator whose spectrum consists of
prime gaps;
\item the formulation of a finite symmetry group acting on prime states,
with explicit Cayley tables and matrix representations;
\item a representation-theoretic decomposition of the prime Hilbert
space into symmetry sectors.
\end{itemize}

This framework places the arithmetic of primes into a setting familiar
from quantum mechanics and noncommutative geometry: operators on Hilbert
spaces, symmetry groups, and spectral decompositions. While the present
paper is primarily structural, it opens a path toward a dynamical theory
of primes governed by operator algebras rather than purely analytic
methods.

% =======================
% CHARACTER TABLE
% =======================

\section{Character Table of the Dihedral Core $D_4$}

For concreteness we include the character table of the dihedral group
$D_4$, which appears as a core symmetry sector of the group $\mathcal G$.

\begin{center}
\begin{tabular}{c|ccccc}
 & $[e]$ & $[r^2]$ & $[r,r^3]$ & $[s,sr^2]$ & $[sr,sr^3]$\\ \hline
$\chi_1$ (trivial) & $1$ & $1$ & $1$ & $1$ & $1$\\
$\chi_2$ & $1$ & $1$ & $1$ & $-1$ & $-1$\\
$\chi_3$ & $1$ & $1$ & $-1$ & $1$ & $-1$\\
$\chi_4$ & $1$ & $1$ & $-1$ & $-1$ & $1$\\
$\chi_5$ (2-dim) & $2$ & $2$ & $0$ & $0$ & $0$
\end{tabular}
\end{center}

Here the conjugacy classes are:
\begin{align*}
[e] &= \{e\},\\
[r^2] &= \{r^2\},\\
[r,r^3] &= \{r,r^3\},\\
[s,sr^2] &= \{s,sr^2\},\\
[sr,sr^3] &= \{sr,sr^3\}.
\end{align*}

The irreducible characters $\chi_i$ define the projectors
\begin{equation}
\Pi_i=\frac{d_i}{|D_4|}\sum_{g\in D_4}\overline{\chi_i(g)}\,U(g),
\end{equation}
which extract the symmetry-adapted components of the prime Dirac field.

% =======================
% POSITIONING & OUTLOOK
% =======================

\section{Outlook}

The group-theoretic Dirac model for primes developed here suggests
several directions for further research. Among them are:
\begin{itemize}
\item the extension from finite symmetry groups to profinite or adelic
symmetry groups acting on primes;
\item the construction of genuinely dynamical operators whose spectra
evolve under symmetry-breaking perturbations;
\item connections with noncommutative geometry, in particular with
spectral triples associated to arithmetic spaces.
\end{itemize}

In this sense the present paper should be viewed as a structural
foundation: it formulates primes and gaps in the language of operators,
symmetries, and spectra, thereby opening the door to a fully fledged
``quantum theory of primes.''
